% TODO make hw stylesheet
\documentclass[10pt]{scrartcl}
\usepackage[a4paper, total={6in, 8in}]{geometry}
\usepackage[usenames,dvipsnames,svgnames]{xcolor}
\usepackage[shortlabels]{enumitem}
\usepackage[framemethod=TikZ]{mdframed}
\usepackage{amsmath,amssymb,amsthm}
\usepackage[colorlinks]{hyperref}
\usepackage{multicol}
\usepackage{tabularx}

\hypersetup{
	colorlinks=true,
	linkcolor=blue,
	filecolor=magenta,      
	urlcolor=magenta,
	pdftitle={MAT 362 HW}
}

\usepackage{microtype}
\usepackage{mathtools}
\usepackage[headsepline]{scrlayer-scrpage}
\usepackage{thmtools}
\usepackage[textsize=footnotesize]{todonotes}

\mdfdefinestyle{mdbluebox}{roundcorner=10pt,innerbottommargin=9pt,
	linecolor=blue,backgroundcolor=TealBlue!5,}
\declaretheoremstyle[headfont=\sffamily\bfseries\color{MidnightBlue},
mdframed={style=mdbluebox},]{thmbluebox}

\mdfdefinestyle{mdbloobox}{frametitlefont=\bfseries,innerbottommargin=8pt,
	nobreak=true,backgroundcolor=TealBlue!5,linecolor=blue,}
\declaretheoremstyle[headfont=\bfseries\color{MidnightBlue},
mdframed={style=mdbloobox},headpunct={\\[3pt]},postheadspace=0pt,]{thmbloobox}

\mdfdefinestyle{mdredbox}{frametitlefont=\bfseries,innerbottommargin=8pt,
	nobreak=true,backgroundcolor=Salmon!5,linecolor=RawSienna,}
\declaretheoremstyle[headfont=\bfseries\color{RawSienna},
mdframed={style=mdredbox},headpunct={\\[3pt]},postheadspace=0pt,]{thmredbox}

\mdfdefinestyle{mdgreenbox}{linecolor=ForestGreen,backgroundcolor=ForestGreen!5,
	linewidth=2pt,rightline=false,leftline=true,topline=false,bottomline=false,}
\declaretheoremstyle[headfont=\bfseries\sffamily\color{ForestGreen!70!black},
mdframed={style=mdgreenbox},headpunct={ --- },]{thmgreenbox}

\mdfdefinestyle{mdorangebox}{linecolor=RedOrange,backgroundcolor=RedOrange!5,
	linewidth=2pt,rightline=false,leftline=true,topline=false,bottomline=false,}
\declaretheoremstyle[headfont=\bfseries\sffamily\color{RedOrange!70!black},
mdframed={style=mdorangebox},headpunct={ --- },]{thmorangebox}

\mdfdefinestyle{mdblackbox}{linecolor=black,backgroundcolor=RedViolet!5!gray!5,
	linewidth=3pt,rightline=false,leftline=true,topline=false,bottomline=false,}
\declaretheoremstyle[mdframed={style=mdblackbox}]{thmblackbox}

\declaretheoremstyle[spaceabove=6pt,spacebelow=6pt]{basehead}

\declaretheorem[style=basehead,name=Problem]{problem}
\declaretheorem[style=basehead,name=Problem,numbered=no]{problem*}
\declaretheorem[style=thmbloobox,name=Setup]{setup} % for config geo
\declaretheorem[style=thmredbox,name=Example,sibling=problem]{example}
\declaretheorem[style=thmredbox,name=$\bigstar$ Problem,sibling=problem]{niceprob}
\declaretheorem[style=thmbluebox,name=Theorem,sibling=problem]{theorem}
\declaretheorem[style=thmbluebox,name=Corollary,sibling=theorem]{corollary}
\declaretheorem[style=thmbluebox,name=Proposition,sibling=theorem]{prop}
\declaretheorem[style=thmbluebox,name=Lemma,numberwithin=problem]{lemma}
\declaretheorem[style=thmbluebox,name=Theorem,numbered=no]{theorem*}
\declaretheorem[style=thmbluebox,name=Lemma,numbered=no]{lemma*}
\declaretheorem[style=thmgreenbox,name=Claim,numberwithin=problem]{claim}
\declaretheorem[style=thmgreenbox,name=Claim,numbered=no]{claim*}
\declaretheorem[style=thmblackbox,name=Remark,numberwithin=problem]{remark}
\declaretheorem[style=thmblackbox,name=Remark,numbered=no]{remark*}
\declaretheorem[style=thmgreenbox,name=Definition,sibling=theorem]{definition}
\declaretheorem[style=thmgreenbox,name=Definition,numbered=no]{definition*}
\declaretheorem[style=thmorangebox,name=Warning,sibling=theorem]{warning}
\declaretheorem[style=thmorangebox,name=Warning,numbered=no]{warning*}
\declaretheorem[style=thmorangebox,name=Note,numbered=no]{note}
\declaretheorem[style=thmblackbox,name=Exercise,sibling=problem]{exercise}
\declaretheorem[style=thmblackbox,name=Exercise,numbered=no]{exercise*}
\declaretheorem[style=thmblackbox,name=Question,sibling=problem]{question}
\declaretheorem[style=thmblackbox,name=Question,numbered=no]{question*}

\addtolength{\textheight}{3.14cm}
\providecommand{\ol}{\overline}
\providecommand{\ul}{\underline}
\providecommand{\eps}{\varepsilon}
\providecommand{\half}{\frac{1}{2}}
\providecommand{\dang}{\measuredangle} %% Directed angle
\providecommand{\CC}{\mathbb C}
\providecommand{\FF}{\mathbb F}
\providecommand{\NN}{\mathbb N}
\providecommand{\QQ}{\mathbb Q}
\providecommand{\RR}{\mathbb R}
\providecommand{\ZZ}{\mathbb Z}
\providecommand{\dg}{^\circ}
\providecommand{\ii}{\item}
\providecommand{\setdiff}{\smallsetminus}
\providecommand{\alert}[1]{{\sffamily\textbf{\textcolor{blue}{#1}}}}
\providecommand{\inv}{^{-1}}
\providecommand{\mb}[1]{\mathbf{#1}}
\providecommand{\dif}{\;d}
\newcommand*{\horzbar}{\rule[.5ex]{2.5ex}{0.5pt}}

\reversemarginpar

\newenvironment{soln}{\begin{proof}[Solution]}{\end{proof}}
\newenvironment{walkthrough}{\noindent \textbf{Walkthrough.}}{\medskip}
\newlist{walk}{enumerate}{3}
\setlist[walk]{label=\bfseries (\alph*)}

\providecommand{\pow}{\operatorname{Pow}}
\providecommand{\vocab}[1]{{\textbf{\textcolor{ForestGreen}{#1}}}}
\providecommand{\lcm}{\operatorname{lcm}}
\providecommand{\abs}[1]{\left|#1\right|}

\usepackage{asymptote}
\begin{asydef}
	defaultpen(fontsize(10pt));
	unitsize(1cm);
	size(8cm); // set a reasonable default
	usepackage("amsmath");
	usepackage("amssymb");
	settings.tex="pdflatex";
	settings.outformat="pdf";
	// Replacement for olympiad+cse5 which is not standard
	import geometry;
	// recalibrate fill and filldraw for conics
	void filldraw(picture pic = currentpicture, conic g, pen fillpen=defaultpen, pen drawpen=defaultpen)
	{ filldraw(pic, (path) g, fillpen, drawpen); }
	void fill(picture pic = currentpicture, conic g, pen p=defaultpen)
	{ filldraw(pic, (path) g, p); }
	// some geometry
	pair foot(pair P, pair A, pair B) { return foot(triangle(A,B,P).VC); }
	pair orthocenter(pair A, pair B, pair C) { return orthocentercenter(A,B,C); }
	pair centroid(pair A, pair B, pair C) { return (A+B+C)/3; }
	// cse5 abbreviations
	path CP(pair P, pair A) { return circle(P, abs(A-P)); }
	path CR(pair P, real r) { return circle(P, r); }
	pair IP(path p, path q) { return intersectionpoints(p,q)[0]; }
	pair OP(path p, path q) { return intersectionpoints(p,q)[1]; }
	path Line(pair A, pair B, real a=0.6, real b=a) { return (a*(A-B)+A)--(b*(B-A)+B); }
	// cse5 more useful functions
	picture CC() {
		picture p=rotate(0)*currentpicture;
		currentpicture.erase();
		return p;
	}
	pair MP(Label s, pair A, pair B = plain.S, pen p = defaultpen) {
		Label L = s;
		L.s = "$"+s.s+"$";
		label(L, A, B, p);
		return A;
	}
	pair Drawing(Label s = "", pair A, pair B = plain.S, pen p = defaultpen) {
		dot(MP(s, A, B, p), p);
		return A;
	}
	path Drawing(path g, pen p = defaultpen, arrowbar ar = None) {
		draw(g, p, ar);
		return g;
	}
\end{asydef}

\usepackage{mathrsfs}

\ihead{\footnotesize MAT 362 HW09}
\ohead{\footnotesize Last compiled \today}

\title{\vspace{-2em}MAT 362 HW09}
\author{\large Aditya Pahuja}
\date{\large Due 04/10}
\begin{document}
\maketitle

\begin{problem*}4.2.3.
	Show that a diffeomorphism $\varphi\colon S\to\ol S$
	is an isometry if and only if the arc length of
	any parametrized curve in $S$ is equal to the arc length
	of the image curve by $\varphi$.
\end{problem*}

\begin{soln}
	todo
\end{soln}

\begin{problem*}4.2.4.
	Use the stereographic projection to show that the sphere
	is locally conformal to a plane.
\end{problem*}

\begin{soln}
	todo
\end{soln}

\begin{problem*}4.2.7.
	Let $V$ and $W$ be ($n$-dimensional) vector spaces
	with inner products denoted by $\langle\bullet,\bullet\rangle$
	and let $F\colon V\to W$ be a linear map.
	Prove that the following are equivalent:
	\begin{enumerate}[(a)]
		\ii $\langle F(v_1), F(v_2)\rangle = \langle v_1,v_2\rangle$
		for all $v_1,v_2\in V$.
		\ii $|F(v)| = |v|$ for all $v\in V$.
		\ii If $\{v_1,\dots,v_n\}$ is an orthonormal basis in $V$,
		then $\{F(v_1),\dots,F(v_n)\}$ is an orthonormal basis in $W$.
		\ii There exists an orthonormal basis $\{v_1,\dots,v_n\}$ in $V$
		such that $\{F(v_1),\dots,F(v_n)\}$ is an orthonormal basis in $W$.
	\end{enumerate}
\end{problem*}

\begin{soln}
	(b)$\implies$(a). The magnitude condition can be written as
	\[\langle F(v),F(v)\rangle = |F(v)|^2 = |v|^2 = \langle v,v\rangle.\]
	Then, by using the linearity of $F$, we get
	\[2\langle F(v),F(w)\rangle
	= |F(v+w)|^2 - |F(v)|^2 - |F(w)|^2
	= |v+w|^2 - |v|^2 - |w|^2
	= 2\langle v,w\rangle\]
	for any $v,w\in V$.\\
	
	(a)$\implies$(c). The preservation of the dot product means that,
	given orthonormal basis $\{v_1,\dots,v_n\}$, we have
	\[\langle F(v_i),F(v_j)\rangle
	= \langle v_i,v_j\rangle = \begin{cases}
		1 & i=j\\
		0 & i\ne j
	\end{cases}\]
	which is exactly what it means for $\{F(v_1),\dots,F(v_n)\}$
	to be an orthonormal basis for $W$.\\
	
	(c)$\implies$(d). All one has to do here is prove the existence
	of an orthonormal basis for $V$, since the result of (c) already tells us
	that the image of the basis will be an orthonormal basis for $W$.
	It is well-known that any finite-dimensional vector space
	with an inner product has an orthonormal basis,
	e.g. by using the Gram-Schmidt algorithm.\\
	
	(d)$\implies$(b). Let $\{v_1,\dots,v_n\}$ be an orthonormal basis of $V$
	such that $\{F(v_1),\dots,F(v_n)\}$ is an orthonormal basis of $W$.
	Let $v = a_1v_1 + \cdots + a_nv_n$ be an arbitrary vector in $V$
	written with its coordinates in the basis $\{v_1,\dots,v_n\}$. Then
	\[|v|^2 = \langle v,v\rangle = \langle a_1v_1 + \cdots + a_nv_n,
	a_1v_1+\cdots+a_nv_n\rangle =\sum_{i=1}^n \sum_{j=1}^n
	\langle a_iv_i, a_jv_j\rangle = \sum_{i=1}^n a_i^2\]
	and similarly,
	\[|F(v)|^2 = \langle F(v),F(v)\rangle =
	\sum_{i=1}^n \sum_{j=1}^n
	\langle a_iF(v_i), a_jF(v_j)\rangle = \sum_{i=1}^n a_i^2,\]
	where in both cases we reduce the sum to $\sum_{i=1}^n a_i^2$
	by using the fact that
	$\langle v_i,v_j\rangle = \langle F(v_i),F(v_j)\rangle = 0$
	if $i\ne j$ and $1$ if $i=j$. Thus (d) implies that
	for any $v\in V$, $|v| = |F(v)|$.\\
	
	Therefore we have shown that (b) implies (a) implies
	(c) implies (d) implies (b), so the four statements are equivalent.
\end{soln}

\begin{problem*}4.2.14.
	We say that a differentiable map $\varphi\colon S_1\to S_2$ preserves angles
	when for every $p\in S_1$ and $v_1,v_2\in T_p(S_1)$, we have
	\[\cos(v_1,v_2) = \cos(d\varphi_p(v_1),d\varphi_p(v_2)).\]
	Prove that $\varphi$ is locally conformal iff it preserves angles.
\end{problem*}

\begin{soln}
	Suppose first that $\varphi$ is locally conformal.
	Then given $p\in S_1$, there exists a neighborhood $V\subseteq S_1$
	of $p$ such that $\varphi\colon V\to\varphi(V)$ is conformal.
	Then, as shown in the book, conformal maps preserve angles, so the equality
	in the problem statement is true.
	
	Conversely, suppose that $\varphi$ preserves angles. We wish to prove that
	for each $p\in S_1$, there is a neighborhood $V\subseteq S_1$ of $p$
	such that the restriction $\varphi\colon V\to \varphi(V)$ is conformal.
	Note that angles being preserved means
	\[\frac{\langle v_1,v_2\rangle}{|v_1||v_2|} =
	\frac{\langle d\varphi_p(v_1),d\varphi_p(v_2)\rangle}
	{|d\varphi_p(v_1)||d\varphi_p(v_2)|},\]
	\todo{finsih}
\end{soln}

\begin{problem*}4.2.17.
	Consider a triangle on the unit sphere whose sides are segments
	of loxodromes and do not contain poles.
	Prove that the sum of the interior angles of such a triangle is $\pi$.
\end{problem*}

\begin{soln}
	todo
\end{soln}

\begin{problem*}4.3.1.
	Show that if $\mb x$ is an orthogonal parametrization,
	that is, $F = 0$, then
	\[K = -\frac{1}{2\sqrt{EG}}\left(\left(\frac{E_v}{\sqrt{EG}}\right)_v +
	\left(\frac{G_u}{\sqrt{EG}}\right)_u\right).\]
\end{problem*}

\begin{soln}
	todo
\end{soln}

\begin{problem*}4.3.2.
	
\end{problem*}

\begin{soln}
	todo
\end{soln}

\begin{problem*}4.3.3.
	
\end{problem*}

\begin{soln}
	todo
\end{soln}

\begin{problem*}4.3.4.
	
\end{problem*}

\begin{soln}
	todo
\end{soln}

\begin{problem*}4.3.8.
	
\end{problem*}

\begin{soln}
	todo
\end{soln}


















\end{document}